%% 摘要

\begin{abstract}
    低空无人机在巡检、测绘、导航和搜索等任务中具有主动运动、多视角观测和连续时空变化等特点。与普通图像问答或视频理解任务不同，无人机智能体不仅需要理解外部地标、场景和目标可见性，还需要理解自身相对位置、视角变化和飞行行为。围绕这一问题，本文开展面向无人机双认知时空推理的多模态大模型能力评测与优化方法研究，将无人机具身推理能力划分为自我状态认知和环境情况认知两条主线，并在图像与视频两种媒介下构建六类评测任务。本文实现了从 AerialVLN 仿真场景扫描、语义点云融合、地标资产复核、行为驱动视频生成到证据感知图像任务生成的四阶段自动化工具链，形成包含 12 个测试场景、512 个有效地标、4096 个图像样本和 2048 个视频样本的测试集。实验部分对闭源模型、开源模型和空间推理微调模型进行了统一评测，并从语义正确性、空间定位、时间定位、跨媒介迁移和双认知平衡等角度进行分析。结果表明，当前多模态大模型已经具备一定无人机场景理解能力，但仍普遍存在语义与证据脱节、自我状态认知弱于环境情况认知、跨媒介稳定性不足等问题。在此基础上，本文进一步实现了 UAV-DualCog-Train 训练集构建方案和基于 Qwen 3.5 4B 的两阶段能力优化方法，包括 LoRA 监督微调、轻量 GRPO 奖励设计、统一训练流程以及多层级链路验证，为后续补充正式训练结果与优化分析奠定了实现基础。
    \keywordslist{无人机, 多模态大模型, 双认知, 时空推理, 测试基准}
\end{abstract}

\begin{engabstract}
    Low-altitude unmanned aerial vehicles operate under active motion, multi-view observation, and continuous spatio-temporal changes. Unlike conventional image question answering or video understanding, a UAV agent must understand not only external landmarks, scenes, and target visibility, but also its own relative position, viewpoint transition, and flight behavior. This thesis studies capability evaluation and optimization of multimodal large models for UAV dual-cognition spatio-temporal reasoning. It formulates embodied UAV reasoning along two dimensions, self-state cognition and environment-state cognition, and instantiates them through six image-based and video-based tasks. A four-stage construction toolchain is implemented, covering AerialVLN scene scanning, semantic point-cloud fusion, landmark asset verification, behavior-driven video generation, and evidence-aware image task generation. The current test set contains 12 scenes, 512 valid landmarks, 4,096 image samples, and 2,048 video samples. Experiments evaluate proprietary models, open-weight models, and spatially specialized fine-tuned models under a unified protocol, and analyze semantic correctness, spatial grounding, temporal grounding, cross-media transfer, and dual-cognition balance. The results show that current multimodal large models have obtained partial UAV scene understanding ability, but still suffer from a gap between answers and evidence, weaker self-state cognition than environment-state cognition, and insufficient cross-media stability. Based on these findings, this thesis further implements UAV-DualCog-Train and a two-stage optimization route on Qwen 3.5 4B, including LoRA-based supervised fine-tuning, lightweight GRPO reward design, a unified training pipeline, and multi-level pipeline validation, which establishes the implementation basis for the formal optimization experiments to be completed next.
    \engkeywordslist{UAV, multimodal large model, dual cognition, spatio-temporal reasoning, benchmark}
\end{engabstract}
